% Slajd 8 – ochrona danych i bezpieczeństwo oraz uruchomienie i utrzymanie poza UMK (pkt 5 wyzwania).
\begin{frame}{Prywatność, bezpieczeństwo i~utrzymanie poza UMK}
  \begin{columns}[T,onlytextwidth]
    \begin{column}{0.5\textwidth}
      \begin{block}{Dane użytkowników i~bezpieczeństwo}
        \footnotesize
        \begin{itemize}
          \item \textbf{Bez konta:} mapa, katalog i~zgłaszanie barier -- anonimowo, bez zapisu adresów IP
                i~identyfikatorów urządzeń.
          \item \textbf{Profil potrzeb to preferencje} (schody, winda, cisza), nigdy diagnoza ani informacja
                o~niepełnosprawności -- tak, jak wymaga wyzwanie.
          \item \textbf{Konto bez hasła:} e-mail + jednorazowy kod (15~min, 5~prób); kody i~tokeny jako SHA-256;
                ciasteczko HttpOnly, SameSite, Secure; usunięcie konta jednym przyciskiem.
          \item \textbf{Serwer:} HTTPS (Caddy, Let's Encrypt), limity żądań (10~kB, zdjęcia 4~MB), treści użytkowników
                escapowane (testy), role sprawdzane przy każdym żądaniu; polityka prywatności w~aplikacji.
        \end{itemize}
      \end{block}
    \end{column}
    \begin{column}{0.5\textwidth}
      \begin{block}{Uruchomienie i~utrzymanie -- propozycja}
        \footnotesize
        \begin{itemize}
          \item \textbf{Operator:} zespół Accessly (spółka lub fundacja) albo partner społeczny; UMK bez
                utrzymywania bazy i~bez dostępu do systemów wewnętrznych.
          \item \textbf{Koszt (szacunek):} VPS w~UE ok.~60--120~zł/mies., domena, poczta transakcyjna;
                asystent AI rozliczany od zapytania (ułamki grosza).
          \item \textbf{Aktualizacje} jednym poleceniem (Ansible); dane miasta co 24~h automatycznie,
                OSM skryptem co miesiąc; kopie zapasowe SQLite.
          \item \textbf{Zgłoszenia i~moderacja} w~panelu administratora (kolejka, przejęcia, poprawki);
                przenośność: Docker na dowolną infrastrukturę, otwarte licencje wymienione w~aplikacji.
        \end{itemize}
      \end{block}
    \end{column}
  \end{columns}

  \note[item]{Pytanie jury „jakie dane zbieracie?”: bez konta – żadnych; z~kontem – e-mail, opcjonalna nazwa, profil potrzeb. Tabela w~polityce prywatności w~aplikacji.}
  \note[item]{RODO: profil potrzeb to świadoma decyzja projektowa – wyzwanie wprost prosi, by nie wymagać informacji o~niepełnosprawności.}
  \note[item]{Kwoty to szacunki na podstawie cenników VPS; nie obiecywać konkretnego dostawcy.}
\end{frame}
